% GENERATED FILE. Edit canonical lessons, then run npm run generate:books.
% canonical-source-hash: fnv1a64:741f341aeefbadcc
% canonical-lessons: ZH-C97-xiaoshi, ZH-W97-shi-time, ZH-C97-meitian, ZH-C97-yihou, ZH-C97-reading, ZH-W97-my-own-line

\chapter{My Day in Characters}
\label{ch:zh-my-day-in-characters}
\hlchaptermodality{97}

\begin{chapteropening}
\textbf{By the end of this chapter:} \emph{I can read a short day in characters and write one line of my own.}

Complete ZH-W97-my-own-line: one line of your own Chinese about where you are, what you want and for how long, chosen rather than copied, with every word you have not written in pinyin.
\end{chapteropening}

\section[\texorpdfstring{\zh{小时} (xiǎoshí)}{小时 (xiǎoshí)}]{\zh{小时} (xiǎoshí) — hour}

\label{lesson:ZH-C97-xiaoshi}



\begin{quote}
A new chapter, and it opens with a word for measuring a day.

\begin{itemize}
\raggedright
  \item \emph{Recall:} say \zh{为什么}, then answer it with \zh{因为} — \emph{wèishénme}, \emph{yīnwèi}
  \item \emph{Recall:} say \emph{wǒ bù shūfu}, I'm not feeling well — \textbf{R1}, the chapter you have finished
  \item \emph{Recall:} say \emph{shíhou}, a time; a moment — \textbf{R4}, recognition at distance
  \item \emph{Recall:} say \emph{xiǎoshíhou}, in childhood — \textbf{R3}, durable
\end{itemize}
\end{quote}

\subsection*{You'll want to know: \zh{小时}}
\textbf{\zh{小时}} — \emph{xiǎoshí} — \textquotedblleft{}hour\textquotedblright{}.

\begin{quote}
\textbf{\zh{一个小时}} \emph{yí ge xiǎoshí} — an hour
\end{quote}

The next lesson puts the second character in your hand. It is the shí of \emph{shíhou} and \emph{xiǎoshíhou}.

\begin{cousinweb}
\textbf{\zh{小时}} (\emph{xiǎoshí}) is built from \zh{小} 'small' and \zh{时} 'time'. Both are inherited from Middle Chinese, \emph{sjewX} and \emph{dzyi} (Guangyun, in Baxter's transcription), \emph{dzyi} glossed 'a season; a time'.
\end{cousinweb}

\begin{grammarlens}[title={What you've built}]
\zh{一} before \zh{个} rises: \emph{yí ge xiǎoshí}, an hour. Before the other tones it falls: \emph{yì tiān}, a day; \emph{yì nián}, a year. Time words count with \zh{一} as often as anything does.
\end{grammarlens}

\subsection*{Your turn}
\begin{itemize}
\raggedright
  \item \emph{Say it:} \emph{xiǎoshí}
  \item \emph{Say it:} \emph{xiǎoshí}, once more
  \item \emph{Say it:} \zh{一个小时} — \emph{yí ge xiǎoshí} — an hour
  \item \emph{Say it:} \zh{一天} — \emph{yì tiān} — a day
  \item \emph{Say it:} \zh{一年} — \emph{yì nián} — a year
  \item \emph{Recall:} say \emph{shíhou}, then \emph{xiǎoshí}, and hear the same shí
\end{itemize}

\begin{tcolorbox}[breakable,colback=teal!4,colframe=teal!35!black,title={Before you move on}]
What does \zh{小时} mean? (\textquotedblleft{}Hour\textquotedblright{}.) Say an hour, and listen to \zh{一}. (\emph{yí ge xiǎoshí}: \zh{一} rises before \zh{个}.)
\end{tcolorbox}
\section[\texorpdfstring{\zh{时} (shí)}{时 (shí)}]{\zh{时} — \zh{日}, then the right half of \zh{对}}

\label{lesson:ZH-W97-shi-time}



\begin{quote}
The word first, then the character inside it.

\begin{itemize}
\raggedright
  \item \emph{Recall:} say \emph{xiǎoshí} — \textbf{R1}, one lesson back, before it decays
\end{itemize}

\begin{quote}
\zh{日} \zh{对}
\end{quote}

\emph{Read it:} each one aloud, with nothing said first — \emph{rì}, \emph{duì}
\end{quote}

\subsection*{Script — \zh{日}, then the right half of \zh{对}}
\hlblockfigure{\detokenize{figures/ZH-W97-shi-time-filmstrip.pdf}}{How \zh{时} is written, stroke by stroke}

\begin{quote}
\zh{时}
\end{quote}

\textbf{Two pieces you own, side by side.}

\textbf{\zh{日}, four strokes}, narrow, on the left: the left side, the top that turns down the right side, the middle bar, the bottom bar.

\textbf{The right half of \zh{对}, three strokes}: the bar, the vertical that hooks to the left at its foot, then the dot.

4 + 3 = seven.

\subsection*{Your turn}
\begin{itemize}
\raggedright
  \item \emph{Write it:} \zh{时}, all seven, with the model in front of you
  \item \emph{Cover:} write it again from memory, and count the pieces, not the strokes
\end{itemize}

\begin{tcolorbox}[breakable,colback=teal!4,colframe=teal!35!black,title={Before you move on}]
How many strokes in \zh{时}? (\textbf{Seven}.) Which character is on the left? (\textbf{\zh{日}}.) Where have you written the right half? (\textbf{At the end of \zh{对}}.)

Source: \href{https://raw.githubusercontent.com/chanind/hanzi-writer-data/68d10a4b21150cae5e1ebbd223eed289cf32d90c/data/\%E6\%97\%B6.json}{Hanzi Writer Data \zh{时}, pinned stroke paths and medians}.
\end{tcolorbox}
\section[měitiān]{měitiān — every day}

\label{lesson:ZH-C97-meitian}



\begin{quote}
Before the new one: read the last lesson's character, then a day with \zh{一} in front of it.

\begin{itemize}
\raggedright
  \item \emph{Read it:} \zh{两个小时}, aloud — \emph{liǎng ge xiǎoshí}, two hours
  \item \emph{Recall:} say \zh{一天} — \emph{yì tiān}, a day, with \zh{一} falling
\end{itemize}
\end{quote}

\subsection*{You'll want to know: měitiān}
\emph{měitiān} — \textquotedblleft{}every day\textquotedblright{}.

\begin{quote}
\emph{wǒ měitiān hē chá} — I drink tea every day.
\end{quote}

Its first character comes later; for now it is a word for the ear.

\begin{cousinweb}
\emph{měitiān} is built from měi 'each' and tiān 'day', the \zh{天} you write. Both are inherited from Middle Chinese, \emph{mwojX} and \emph{then} (Guangyun, in Baxter's transcription); the Guangyun glosses \emph{mwojX} 'frequently', and 'each' is assumed from that, not quoted.
\end{cousinweb}

\begin{grammarlens}[title={What you've built}]
A time word goes early, after the person and before the verb: \zh{我}, \emph{měitiān}, \emph{hē chá}. \emph{yì tiān} is one day; \emph{měitiān} is every day.
\end{grammarlens}

\subsection*{Your turn}
\begin{itemize}
\raggedright
  \item \emph{Say it:} \emph{měitiān}
  \item \emph{Say it:} \emph{měitiān}, once more
  \item \emph{Say it:} \emph{wǒ měitiān hē chá} — I drink tea every day.
  \item \emph{Recall:} say \emph{yì tiān}, then \emph{měitiān}
\end{itemize}

\begin{tcolorbox}[breakable,colback=teal!4,colframe=teal!35!black,title={Before you move on}]
What does \emph{měitiān} mean? (\textquotedblleft{}Every day\textquotedblright{}.) Where does it go? (\textbf{After the person, before the verb}.)
\end{tcolorbox}
\section[yǐhòu]{yǐhòu — after; afterwards}

\label{lesson:ZH-C97-yihou}



\begin{quote}
Before the new one: last lesson's word, then two words for school and after.

\begin{itemize}
\raggedright
  \item \emph{Recall:} say \emph{měitiān}, every day — \textbf{R1}, one lesson back
  \item \emph{Recall:} say \zh{下课} — \emph{xiàkè}, to finish class
  \item \emph{Recall:} say \emph{hòu}, after; behind — \textbf{R2}
\end{itemize}
\end{quote}

\subsection*{You'll want to know: yǐhòu}
\emph{yǐhòu} — \textquotedblleft{}after; afterwards\textquotedblright{}.

\begin{quote}
\emph{xiàkè yǐhòu, wǒ qù chī fàn} — After class I go and eat.
\end{quote}

\begin{quote}
\emph{yí ge xiǎoshí yǐhòu} — an hour later
\end{quote}

Its characters come later; for now it is a word for the ear.

\begin{cousinweb}
\emph{yǐhòu} is built from yǐ, the yǐ of \emph{kěyǐ}, and hòu 'after', the \emph{hòu} you know. Both are inherited from Middle Chinese, \emph{yiX} and \emph{huwX} (Guangyun, in Baxter's transcription), \emph{huwX} glossed 'after, in order'.
\end{cousinweb}

\begin{grammarlens}[title={What you've built}]
\emph{yǐhòu} comes after the thing it follows: \emph{xiàkè yǐhòu}, after class; \emph{yí ge xiǎoshí yǐhòu}, an hour later.
\end{grammarlens}

\subsection*{Your turn}
\begin{itemize}
\raggedright
  \item \emph{Say it:} \emph{yǐhòu}
  \item \emph{Say it:} \emph{yǐhòu}, once more
  \item \emph{Say it:} \emph{xiàkè yǐhòu, wǒ qù chī fàn} — After class I go and eat.
  \item \emph{Say it:} \emph{yí ge xiǎoshí yǐhòu} — an hour later
  \item \emph{Recall:} say \emph{hòu}, then \emph{yǐhòu}
\end{itemize}

\begin{tcolorbox}[breakable,colback=teal!4,colframe=teal!35!black,title={Before you move on}]
What does \emph{yǐhòu} mean? (\textquotedblleft{}After; afterwards\textquotedblright{}.) Where does it go? (\textbf{After the thing it follows: \emph{xiàkè yǐhòu}}.)
\end{tcolorbox}
\section[(reading)]{(reading) — a Sunday at home}

\label{lesson:ZH-C97-reading}



\begin{quote}
Fifty-seven characters in a row, and every one is a character you write.

\begin{itemize}
\raggedright
  \item \emph{Recall:} say \emph{xiàkè yǐhòu}, after class — \textbf{R1}, one lesson back
  \item \emph{Recall:} say \emph{jiàoshì}, a classroom — \textbf{R3}, durable
\end{itemize}

\begin{quote}
\zh{上课} \zh{下课}
\end{quote}

\emph{Read it:} both aloud, with nothing said first — \emph{shàngkè}, \emph{xiàkè}
\end{quote}

\begin{tcolorbox}[breakable,skin=enhanced,colback=orange!6,colframe=orange!45!black,boxrule=0.5pt,arc=1mm,left=6pt,right=6pt,top=4pt,bottom=4pt,fonttitle=\bfseries,title={Read this}]
\begin{quote}
\zh{今天是星期天} \zh{我和妈妈都在家} \zh{妈妈做饭} \zh{我看书} \zh{我想喝水} \zh{也想吃饭} \zh{妈妈说} \zh{饭太多了} \zh{我说} \zh{太好了} \zh{下午两点} \zh{我们去买书} \zh{因为我很喜欢看书}
\end{quote}

\emph{Read it:} the passage through once without stopping; do not translate as you go

\emph{Read it:} the passage again, and this time notice where the day turns: the line with the time in it
\end{tcolorbox}

\begin{tcolorbox}[breakable,skin=enhanced,colback=orange!6,colframe=orange!45!black,boxrule=0.5pt,arc=1mm,left=6pt,right=6pt,top=4pt,bottom=4pt,fonttitle=\bfseries,title={Read this}]
\begin{quote}
\emph{jīntiān shì xīngqītiān} \emph{wǒ hé māma dōu zài jiā} \emph{māma zuò fàn} \emph{wǒ kàn shū} \emph{wǒ xiǎng hē shuǐ} \emph{yě xiǎng chī fàn} \emph{māma shuō} \emph{fàn tài duō le} \emph{wǒ shuō} \emph{tài hǎo le} \emph{xiàwǔ liǎng diǎn} \emph{wǒmen qù mǎi shū} \emph{yīnwèi wǒ hěn xǐhuan kàn shū}
\end{quote}

One line per clause, and no punctuation, as this book prints Chinese: the line breaks carry the pauses. \zh{饭太多了} is \zh{太} … \zh{了} as a complaint, too much rice; \zh{太好了} a line later is the same frame as delight.
\end{tcolorbox}

\subsection*{Your turn}
\begin{itemize}
\raggedright
  \item \emph{Say it:} the passage aloud, once, without stopping
\end{itemize}

\emph{Twice through:}

\begin{tcolorbox}[breakable,colback=teal!4,colframe=teal!35!black,title={Before you move on}]
Who is at home? (\textbf{\zh{我和妈妈}}: the writer and Mum.) Why do they go to buy books? (\textbf{\zh{因为我很喜欢看书}}.) Which line is a complaint, and which a delight? (\textbf{\zh{饭太多了}, then \zh{太好了}}.)
\end{tcolorbox}
\section[(my own line)]{My own line — the first Chinese with no model}

\label{lesson:ZH-W97-my-own-line}



\begin{quote}
Everything you have written so far had something to reach for: a model to trace, a word somebody said, a sentence printed above the space.

This does not. You choose the words, from a supply that is closed: what this book has given you.

\begin{itemize}
\raggedright
  \item \emph{Recall:} say \zh{两个小时} — \emph{liǎng ge xiǎoshí}, two hours
\end{itemize}
\end{quote}

\subsection*{Writing — controlled composition}
Write \textbf{one line} of your own about your day. Three parts, in your own words: where you are, with \zh{在}; what you want, with \zh{想要} or \zh{要}; and for how long, with \zh{一个小时} or \zh{两个小时}. If you want a second thing, \zh{也} brings it in, in front of its verb.

Some words you will want have characters you have not written yet: \emph{měitiān}, every day; \emph{yǐhòu}, after. Write those in pinyin, inside the line, with their tone marks. This task allows that, and it allows nothing else: a character you can read but have not learned to write is not one you produce, so it goes in pinyin too.

Three checks belong to the hand. Write each character in the stroke order you learned: \zh{时} is \zh{日} first, narrow, then the hooked vertical and the dot. Give every character \textbf{the same size of square}, two-stroke \zh{了} as much as ten-stroke \zh{都}. And leave a space between clauses: no Chinese comma or full stop, as this book prints them.

Write the line. Then read it back, and check that every pinyin word has its tone marks before you check anything else.

\begin{tcolorbox}[breakable,colback=teal!4,colframe=teal!35!black,title={Before you move on}]
Which words did you write in pinyin? (\textbf{Any whose characters you have not written: \emph{měitiān}, \emph{yǐhòu}}.)

Which piece of \zh{时} comes first? (\textbf{\zh{日}}, narrow, on the left.)

And what did you choose here that the earlier lessons chose for you? (\textbf{The words.})
\end{tcolorbox}
